\chapter{Packaging the Pilot for a Research Collaboration Meeting (Day 14)}
\label{chap:day14}

\section{Motivation}

Chapters~\ref{chap:day01}--\ref{chap:day13} document thirteen days of building, testing, and
verifying a six-metric evaluation pipeline against a vision-language model. None of those
chapters' artifacts are written for the specific occasion this pilot was always aiming toward:
a meeting with Professor Mustafa Abdallah, the collaborator whose six-metric framework this
pilot adapts, who was not present for any of the preceding thirteen days and needs a small set
of focused materials rather than a thirteen-chapter report. Day 14's task, per the two-week
pilot plan, was to finalize that material: a one-page pitch, a curated set of sample
visualizations, a task list for the full study framed as decisions rather than to-dos, and a
five-minute spoken explanation covering what was tested, what worked, what failed, and where
Professor Abdallah's expertise is specifically needed. This chapter is not a fourteenth unit of
evaluation work; it is the closing act of packaging thirteen days of it for the person the work
was done for.

\section{Artifacts Produced}

Four documents were produced on Day 14, each serving a different moment in the meeting rather
than duplicating the same content four times. \texttt{report/one\_pager\_summary.md} is the
single page meant to be read before or during the meeting: the question, the headline answer,
the six metrics' headline numbers, the findings that matter more than those numbers, and the
points where guidance is wanted. \texttt{report/talking\_points.md} is the five-minute spoken
structure built to walk through that same material out loud --- the question (30 seconds), what
was tested (one minute), what worked (1.5 minutes), what failed or is more interesting than the
headline numbers (1.5 minutes), and where guidance is wanted (30 seconds) --- speaking notes
rather than a script. \texttt{report/key\_visualizations.md} curates six images out of the
60-plus saved across this pilot's thirteen days, each selected as the clearest single-image
evidence of a specific finding rather than for visual appeal. \texttt{report/full\_study\_task\_list.md}
converts six open items into six explicit decisions, each naming the chapter whose evidence
motivates it, so the meeting's task-list discussion starts from pilot evidence rather than a
generic future-work list.

Every number in all four documents traces to a results file or a chapter already in this
report; none was invented or estimated for the meeting package. This is the same sourcing
discipline Chapter~\ref{chap:day12}'s pilot-report memo held under its own 14-day timebox,
inherited here rather than relaxed for what is, in volume, the smallest of this pilot's written
artifacts.

\section{The Headline Answer}
\label{sec:day14-headline}

Chapter~\ref{chap:day01} opened this report by establishing that the pilot was executable at
all; underneath that operational question sat the substantive one this entire pilot exists to
answer: \textbf{does Professor Abdallah's six-metric XAI evaluation framework --- Descriptive
Accuracy, Sparsity, Stability, Efficiency, Robustness, and Completeness, built and validated on
tabular intrusion-detection DNNs --- transfer to a vision-language model making visual
judgments?} Thirteen chapters of evidence now stand behind a direct answer.

\textbf{Yes, with two real, quantified adaptations --- not zero changes.} All six metrics ran
end-to-end against Florence-2-base-ft (231M parameters, a single RTX 3070) on 163
construction-safety images sampled from \texttt{LouisChen15/ConstructionSite}'s 3,004-image
test split. Neither adaptation was optional or cosmetic; each was a response to a concrete
failure the framework's own measurements surfaced, not a precaution taken in advance of trouble.

\textbf{The first adaptation} is Chapter~\ref{chap:day03}'s grounding-proxy reframing.
Florence-2 has no native free-form visual-question-answering token, so each safety rule had to
be reframed as a grounding query against the model's actual task vocabulary
(\texttt{<OPEN\_VOCABULARY\_DETECTION>}) rather than posed as a literal yes/no question. This
adaptation sits at the model-interface level, not the metric level: none of the six metrics had
to be redefined to accommodate it, but every one of them depends on the resulting proxy answer
being a meaningful judgment in the first place.

\textbf{The second adaptation} is the person-relative redesign, also from
Chapter~\ref{chap:day03}, found only after the first version of the grounding proxy was tested
and found wanting: a scene-level existence check achieved 3.0\% sensitivity, a near-constant
``compliant'' classifier. Redesigning the proxy to check coverage per detected worker, mirroring
the dataset's own per-worker labeling semantics, produced a \textbf{9$\times$ sensitivity gain
(27.0\%)}. This is the more consequential of the two adaptations: Descriptive Accuracy and
Bounded Completeness, the framework's own metrics, measure whether an \emph{explanation} is
faithful to a judgment --- they do not, by themselves, measure whether the underlying
judgment-generation method is discriminating at all. A near-constant classifier would have left
those metrics almost nothing real to test, and the framework's own machinery would not have
flagged that on its own; it took a separate sensitivity/specificity check, run before any
explanation method was built on the proxy, to catch it.

With both adaptations in place, all six metrics produced genuine, interpretable signal: 36.2\%
descriptive accuracy, 0.108 top-5 visual sparsity, 77.5\% stability, 80.7\% Level-1 robustness,
43.6\% bounded completeness, and a confirmed-trivial $\sim$0.5 GPU-hours per 1,000 samples for
the full pipeline. None of the six metrics had to be abandoned or silently redefined into
something the framework's original authors would not recognize. The honest reading of these
thirteen chapters is that the framework's \emph{structure} transferred intact across the gap
from tabular intrusion-detection features to image-and-text VLM inputs; what did not transfer
for free was the assumption, implicit in applying any of the six metrics, that the system being
evaluated already produces a discriminating judgment to explain. That assumption needed to be
checked and, on first attempt, repaired --- exactly the kind of adaptation a feasibility pilot
exists to surface before a larger study is built on an unverified premise.

\section{Curated Visualizations}

\texttt{report/key\_visualizations.md} curates six images out of the more than sixty saved
across Chapters~\ref{chap:day03}--\ref{chap:day11}, selected because each is the clearest
single-image evidence for one specific finding, not for being the most visually striking option
available. Each is identified here by the chapter that produced it and the finding it
illustrates; none is re-embedded in this chapter, since each already appears, in full, at the
point in this report where it was first introduced and explained.

\begin{enumerate}
\item \textbf{The headline rollup} (\texttt{summary/metric\_summary\_by\_class.png}, first
shown as Figure~\ref{fig:day11-summary}). Chapter~\ref{chap:day11}'s capstone chart: all five
per-sample metrics, broken down by hazard class, on their own native scales. This is the single
chart that the rest of the curated set exists to explain --- everything else on this list
accounts for some feature of why this chart looks the way it does.

\item \textbf{The grounding-proxy adaptation, working} (\texttt{baseline/0000007\_rule\_1.png},
first shown as Figure~\ref{fig:day03-rule1}). Chapter~\ref{chap:day03}'s worked example of the
redesigned person-relative proxy in action: a hard-hat query correctly localized on a worker's
head, the underlying mechanism behind the 3.0\%-to-27.0\% sensitivity gain that is this pilot's
second headline adaptation.

\item \textbf{Descriptive accuracy on a real sample}
(\texttt{descriptive\_accuracy/0000037\_rule\_1\_flip.png}). The clearest single-image
illustration, among Chapter~\ref{chap:day05}'s ten saved flip overlays, of masking the top-ranked
region changing the model's answer --- descriptive accuracy behaving exactly as the metric is
designed to detect.

\item \textbf{Sparsity, focused versus scattered.}
\texttt{sparsity/0000023\_rule\_1\_hard\_hat.png} and
\texttt{sparsity/0000079\_rule\_3\_guardrail.png}, first shown together as
Figure~\ref{fig:day06-worked}. Chapter~\ref{chap:day06}'s side-by-side pair makes the
chapter's size-confound finding visible without requiring the correlation number: a small,
sharp hotspot on a distant hard hat against a broadly scattered heatmap for a frame-spanning
guardrail, two physically different targets rather than two differently-explained judgments.

\item \textbf{A stable answer hiding a drifted explanation}
(\texttt{robustness/verify/0000034\_before.png} and
\texttt{robustness/verify/0000034\_after\_occlude.png}, first shown together as
Figure~\ref{fig:day09-finding2}). The single best image in this pilot for
Chapter~\ref{chap:day09}'s centerpiece finding: the same final answer before and after
occlusion, while the model's evidence box quietly relocates to something unrelated
(\texttt{object\_box\_iou=0.005}) --- the case that motivates reporting answer-level and
explanation-level robustness as two never-blended numbers.

\item \textbf{Compute is not the obstacle} (\texttt{efficiency/per\_sample\_cost.png}, first
shown as Figure~\ref{fig:day08-worked}). The per-sample timing breakdown behind
Chapter~\ref{chap:day08}'s $\sim$0.5-GPU-hours-at-$n$=1,000 figure, the evidence that efficiency
is the one metric in this pilot whose answer is unqualified good news for scaling up.
\end{enumerate}

\section{Six Decisions for Professor Abdallah}
\label{sec:day14-decisions}

\texttt{report/full\_study\_task\_list.md} converts six items this pilot deliberately scoped
down or left open into six explicit decisions, each framed as something requiring Professor
Abdallah's judgment before it is worth committing engineering time to, and each tied to the
specific chapter whose evidence motivates it. None of the six is something the pilot attempted
and failed at; each is a fork in the road the pilot reached and stopped at, by design, rather
than a defect to be silently fixed.

\textbf{1. Attribution method (Chapter~\ref{chap:day06}).} Is decoder$\to$encoder
cross-attention an acceptable stand-in for LRP or DeepLIFT in the framework's own terms, or does
a transformer-native attribution method need to be added before the framework's attribution
component can be considered fully ported? Chapter~\ref{chap:day06} verified directly that
classic attention rollout's single-modality assumption does not hold for Florence-2's fused
image+text encoder, and substituted cross-attention only after that verification, but the
substitution's acceptability as formal evidence for the framework's purposes is a judgment call
this pilot is not positioned to make on its own.

\textbf{2. Statistical rigor at small $n$, especially \texttt{struck\_by\_risk}
(Chapters~\ref{chap:day02}, \ref{chap:day11}).} Chapter~\ref{chap:day02} found that the entire
3,004-image test split contains only 13 \texttt{struck\_by\_risk} examples, capping that class's
sample size well below the other three classes' 50 each; Chapter~\ref{chap:day11} flagged every
\texttt{struck\_by\_risk} number in its headline table as a directional signal rather than a
stable estimate for exactly this reason. What minimum $n$ per class should be required before a
metric is reported as a stable estimate, and what test or interval should accompany each
headline number going forward, is a decision this pilot defers rather than resolves.

\textbf{3. Region-ranking heuristic (Chapters~\ref{chap:day04}, \ref{chap:day10}).} Should
candidate regions be ranked by whether they match the rule's queried object class before
falling back to raw area, removing the worker-versus-PPE-object confound that
Chapter~\ref{chap:day04} flagged as a foreseen limitation at the moment
\texttt{standardize\_regions} was designed and that Chapter~\ref{chap:day10} later measured
directly: masking the worker flips the answer in only 37.5\% of cases against 49.4\% for the
actual safety object, an 11.9 percentage-point gap traced to one specific, replaceable ranking
function. This is the single highest-leverage, most fixable item on this list, and fixing it
plausibly improves three metrics' weakest results simultaneously for PPE classes --- but it is
listed as a decision, not an authorized fix, because it changes the regions every later
masking-based chapter in this pilot depends on.

\textbf{4. Full-dataset scaling versus the \texttt{struck\_by\_risk} class imbalance
(Chapter~\ref{chap:day08}).} Is scaling to the full 3,004-image test split the right next step,
or is correcting the 13-example \texttt{struck\_by\_risk} shortage a higher-value use of the
same engineering effort? Chapter~\ref{chap:day08} confirmed that compute is not the deciding
factor either way: the full six-metric pipeline costs roughly half a GPU-hour at $n$=1,000 on a
single consumer-grade RTX 3070, so this decision should turn on what scaling or rebalancing
strategy best serves the eventual paper's evidentiary needs, not on what a single RTX 3070 can
afford to run.

\textbf{5. A real Level-3 robustness study (Chapter~\ref{chap:day09}).} What scope and budget
should a properly powered adversarial-patch study receive? Chapter~\ref{chap:day09}'s 20-image
stretch test had only 4 candidates for which a flip to ``compliant'' was even a meaningful
outcome, and those four cases showed no clean placement-to-flip correlation --- a well-placed
patch failed to fool the model, a badly-placed one succeeded. That result is a plausibility
signal worth a properly powered follow-up, explicitly not a result strong enough to support a
claim about Florence-2's vulnerability to adversarial patches on its own.

\textbf{6. \texttt{rule\_2}'s prompt-phrasing ambiguity (Chapter~\ref{chap:day09}, Finding~3).}
Should \texttt{rule\_2}'s two phrasings, ``safety harness'' and ``fall-protection lanyard,'' be
treated as two distinct physical concepts rather than wording variants of one concept, before
\texttt{rule\_2}'s robustness numbers are trusted in a larger study? Chapter~\ref{chap:day09}
found a 50.0\% Level-2 answer-change rate for \texttt{rule\_2} against 10.2\% for \texttt{rule\_3}
under the same rewording test, and traced the gap to the two \texttt{rule\_2} phrasings
plausibly grounding two different objects on the same scene, not to any instability in how
Florence-2 handles a reworded question.

\section{Standard vs. Non-Standard Choices}

\textbf{Standard}: assembling a stakeholder-facing summary package --- a one-pager, talking
points, and a curated set of visuals --- at the close of a feasibility pilot, ahead of a
collaboration meeting, is itself an entirely ordinary practice. Nothing about producing four
short documents to support a five-minute spoken explanation departs from how any applied
research pilot would normally close out before reporting to a collaborator or sponsor.

\textbf{Non-standard}, and the choice worth naming explicitly: framing every open item in the
full study task list as an explicit decision that names the specific pilot evidence behind it,
rather than as a generic ``future work'' list compiled after the fact. A generic future-work
list would have let each of the six items in Section~\ref{sec:day14-decisions} read as an
unexamined gap; instead, each is presented as a fork already reached, with the chapter and
finding that produced it cited directly, so that Professor Abdallah's judgment is brought to
bear on a decision with evidence already attached rather than on an open-ended request for
direction. This mirrors, at the level of the whole pilot's closing package, the same discipline
Chapter~\ref{chap:day12}'s pilot-report memo adopted at the level of a single document: every
number traces to a specific source, and every open question traces to a specific chapter's
finding, rather than either being supplied from memory or general impression.

\section{Closing Synthesis}

Chapter~\ref{chap:day01} began this report by removing one assumption --- that the pilot was
executable at all --- so that every later day's plan would rest on verified ground rather than
on faith. Thirteen chapters later, the report closes by answering the question that operational
verification was always in service of: \textbf{does Professor Abdallah's six-metric XAI
evaluation framework, built and validated on tabular intrusion-detection DNNs, transfer to a
vision-language model making visual judgments on construction-safety imagery?} Section~\ref{sec:day14-headline}
gave that answer directly --- yes, with two real, quantified adaptations, not zero changes ---
and it is worth restating here, at the report's close, why that answer is the right level of
confidence to report rather than either a more triumphant or a more hedged one.

It would overstate the evidence to claim the framework transferred unmodified: Chapter~\ref{chap:day03}'s
person-relative redesign was a 9$\times$ change in the underlying proxy's discriminative power,
found only because the adoption process included a direct sensitivity check the six metrics
themselves do not perform. It would equally understate the evidence to treat that adaptation as
a sign the framework does not generalize: every one of the six metrics, once given a
discriminating judgment to explain, ran to completion and produced real, interpretable,
traceable signal --- Chapter~\ref{chap:day05}'s non-monotonicity anomaly, Chapter~\ref{chap:day06}'s
size confound, Chapter~\ref{chap:day09}'s answer-level-versus-explanation-level split, and
Chapter~\ref{chap:day10}'s area-ranking finding are exactly the kind of findings a working
evaluation framework produces when pointed at a real system, not symptoms of a framework
failing to fit its new target. The thirteen chapters behind this answer were written
contemporaneously, day by day, as Chapter~\ref{chap:day12} recorded, and the cross-cutting
patterns connecting them --- the fallback-rerouting mechanism recurring across
Chapters~\ref{chap:day05}, \ref{chap:day09}, and \ref{chap:day10}; the area-ranking confound
recurring across Chapters~\ref{chap:day04}, \ref{chap:day06}, \ref{chap:day07}, and
\ref{chap:day10} --- were connected only afterward, on a rereading of an already-written
record, not retrofitted to manufacture a tidier story than the evidence supports. That is the
discipline this closing synthesis relies on: the six-metric framework transfers to a
vision-language model, and the report's own day-by-day record is the reason that claim can be
made with the confidence it is made here.

\section{Implications for the Paper}

This chapter is the natural home for the eventual paper's Conclusion and Future Work section,
and the two pieces of this chapter map onto that section almost directly. Section~\ref{sec:day14-headline}'s
headline answer --- transfer succeeds, with two quantified, named adaptations --- is the
Conclusion: a single paragraph the paper's abstract and introduction can both be built around,
with Chapters~\ref{chap:day01}--\ref{chap:day13} supplying the supporting evidence in full
behind it. Section~\ref{sec:day14-decisions}'s six decisions are the Future Work section's
candidate bullets almost without rewriting: each already names the specific chapter and finding
that motivates it, which is precisely the structure a Future Work section should have ---
specific, evidenced, and prioritized, rather than a closing paragraph of generic caveats.

Two of the six decisions deserve to be highlighted as priorities if the paper's Future Work
section needs to rank rather than simply list them. The region-ranking heuristic (decision~3) is
the single highest-leverage fix identified anywhere in this pilot, named as a foreseen
limitation as early as Chapter~\ref{chap:day04} and measured directly by Chapter~\ref{chap:day10}:
it requires no new model calls, only a re-ranking of boxes Florence-2 has already returned, and
Chapter~\ref{chap:day10} already supplies a falsifiable prediction for what a successful fix
should produce. The attribution-method question (decision~1) is the most consequential for how
the paper positions its own methodological contribution: whether decoder$\to$encoder
cross-attention is accepted as sufficient evidence for a fused-modality vision-language model,
or whether the framework requires a transformer-native method before its attribution component
can be considered fully ported, determines how strongly the paper can claim the full framework,
rather than five of its six components, has been validated on a VLM. Both decisions, and the
four alongside them, are now in Professor Abdallah's hands with this pilot's evidence already
attached --- which is the entire purpose this closing chapter, and the meeting it was written
for, was built to serve.
